% Part: second-order-logic
% Chapter: metatheory
% Section: second-order-arithmetic

\documentclass[../../../include/open-logic-section]{subfiles}

\begin{document}

\olfileid{sol}{met}{spa}

\olsection{દ્વિતીય-ક્રમ અંકગણિત}

યાદ કરો કે પેયાનો અંકગણિતના સિદ્ધાંત~$\Th{PA}$માં~$\Th{Q}$નાં આઠ
સ્વયંસિદ્ધિઓ,
\begin{align*}
& \lforall[x][\eq/[x'][\Obj 0]]\\
& \lforall[x][\lforall[y][(\eq[x'][y'] \lif \eq[x][y])]]\\
& \lforall[x][(\eq[x][\Obj 0] \lor \lexists[y][\eq[x][y']])]\\
& \lforall[x][\eq[(x + \Obj 0)][x]]\\
& \lforall[x][\lforall[y][\eq[(x + y')][(x + y)']]]\\
& \lforall[x][\eq[(x \times \Obj 0)][\Obj 0]]\\
& \lforall[x][\lforall[y][\eq[(x \times y')][((x \times y) + x)]]]\\
& \lforall[x][\lforall[y][(x < y \liff \lexists[z][\eq[(z' + x)][y]])]]\\
\intertext{અને આ સ્વરૂપનાં બધાં વાક્યો સામેલ છે:}
& (!A(\Obj 0) \land \lforall[x][(!A(x) \lif !A(x'))]) \lif \lforall[x][!A(x)].
\end{align*}
છેલ્લું ``પ્રરૂપ'' છે, એટલે અંકગણિતની ભાષાના દરેક
!!{formula}~$!A(x)$ માટે એક, એમ અનંત સંખ્યામાં !!{sentence}s
ઉત્પન્ન કરતું માળખું. આ પ્રરૂપને (પ્રથમ-ક્રમ) \emph{આગમનનું
સ્વયંસિદ્ધિ-પ્રરૂપ} કહીએ છીએ.
\emph{દ્વિતીય-ક્રમ} પેયાનો અંકગણિત~$\Th{PA^2}$માં આગમનને એક જ
વાક્યથી વ્યક્ત કરી શકાય છે. $\Th{PA^2}$ ઉપરની પ્રથમ આઠ સ્વયંસિદ્ધિઓ
અને (દ્વિતીય-ક્રમ) \emph{આગમનના સ્વયંસિદ્ધિ}થી બને છે:
\[
\lforall[X][((X(\Obj 0) \land \lforall[x][(X(x) \lif X(x'))]) \lif \lforall[x][X(x)])].
\]
તે કહે છે: !!{domain}નો ઉપગણ~$X$ જો~$\Assign{\Obj{0}}{M}$ને સમાવે
અને કોઈપણ~$x \in \Domain{M}$ સાથે~$\Assign{\prime}{M}(x)$ને પણ
સમાવે (એટલે કે તે ``અનુગામી હેઠળ બંધ'' હોય), તો તે આખું
!!{domain} સમાવે (એટલે $X =
\Domain{M})$.

આગમનનું સ્વયંસિદ્ધિ ખાતરી આપે છે કે તેને સંતોષતી કોઈપણ !!{structure}ના
$\Domain{M}$માં સ્વયંસિદ્ધિઓથી હોવા જરૂરી હોય એવા !!{element}s જ છે,
એટલે $n \in \Nat$ માટે~$\num{n}$નાં મૂલ્યો. $\Th{PA^2}$ના
નિદર્શમાં કોઈ અપ્રમાણભૂત સંખ્યાઓ નથી.

\begin{thm}
\ollabel{thm:sol-pa-standard}
જો $\Sat{M}{\Th{PA^2}}$ હોય, તો $\Domain{M} =
\Setabs{\Value{\num{n}}{M}}{n \in \Nat}$.
\end{thm}

\begin{proof}
$N = \Setabs{\Value{\num{n}}{M}}{n \in \Nat}$ લો અને
$\Sat{M}{\Th{PA^2}}$ માનો. દરેક $n \in \Nat$ માટે
$\Value{\num{n}}{M} \in \Domain{M}$ છે જ; તેથી
$N \subseteq \Domain{M}$.

હવે બીજી દિશાનો સમાવેશ બતાવીએ. $s(X) = N$ એવી ચલ-નિયુક્તિ~$s$
વિચારો. ધારણા મુજબ,
\begin{align*}
\Sat{M}{\lforall[X][((X(\Obj 0) \land \lforall[x][(X(x)
      \lif X(x'))]) \lif \lforall[x][X(x)])]}, & \text{ તેથી}\\
\Sat{M}{(X(\Obj 0) \land \lforall[x][(X(x) \lif X(x'))]) \lif
  \lforall[x][X(x)]}[s]. &
\end{align*}
આ શરતી વિધાનનો પૂર્વપક્ષ વિચારો. $\Value{\Obj{0}}{M} \in
N$ છે; તેથી $\Sat{M}{X(\Obj{0})}[s]$. બીજો સંયોજક,
$\lforall[x][(X(x) \lif X(x'))]$, પણ સંતોષાય છે. કારણ કે
$x \in N$ હોય એમ માનો. $N$ની વ્યાખ્યા મુજબ કોઈ~$n$ માટે
$x = \Value{\num{n}}{M}$. તેથી
$\Assign{\prime}{M}(x) = \Value{\num{n+1}}{M} \in N$. આમ,
$\Assign{\prime}{M}(x) \in N$.

આથી $\Sat{M}{X(\Obj 0) \land \lforall[x][(X(x) \lif X(x'))]}[s]$.
પરિણામે $\Sat{M}{\lforall[x][X(x)]}[s]$. પરંતુ તેનો અર્થ એ છે કે
દરેક $x \in \Domain{M}$ માટે $x \in s(X) = N$. તેથી
$\Domain{M} \subseteq N$.
\end{proof}

\begin{cor}
\ollabel{cor:sol-pa-categorical}
$\Th{PA^2}$નાં કોઈપણ બે નિદર્શો એકરૂપ છે.
\end{cor}

\begin{proof}
\olref{thm:sol-pa-standard} મુજબ $\Th{PA^2}$ના કોઈપણ નિદર્શનું
વ્યાપકક્ષેત્ર $\Value{\num{n}}{M}$નાં મૂલ્યોથી પૂરું થાય છે. આવું
દરેક નિદર્શ~$\Th{Q}$નું નિદર્શ પણ છે.
\olref[mod][mar][stm]{prop:thq-standard} મુજબ આવું દરેક નિદર્શ
પ્રમાણભૂત છે, એટલે~$\Struct{N}$ને એકરૂપ છે.
\end{proof}

ઉપર આપણે~$\Th{PA^2}$ને પ્રથમ આઠ અંકગણિતીય સ્વયંસિદ્ધિઓ તથા
દ્વિતીય-ક્રમ આગમનનું સ્વયંસિદ્ધિ ધરાવતા સિદ્ધાંત તરીકે વ્યાખ્યાયિત
કર્યો. હકીકતમાં દ્વિતીય-ક્રમ તર્કશાસ્ત્રની અભિવ્યક્તિની શક્તિને
કારણે દ્વિતીય-ક્રમ પેયાનો અંકગણિત માટે અંકગણિતીય સ્વયંસિદ્ધિઓમાંથી
માત્ર \emph{પ્રથમ બે} અને આગમન પૂરતાં છે.

\begin{prop}\ollabel{prop:sol-pa-definable}
$\Th{PA^{2\dagger}}$ એ પ્રથમ બે અંકગણિતીય સ્વયંસિદ્ધિઓ (અનુગામી
વિશેની સ્વયંસિદ્ધિઓ) અને દ્વિતીય-ક્રમ આગમનનું સ્વયંસિદ્ધિ ધરાવતો
દ્વિતીય-ક્રમ સિદ્ધાંત હોય. ત્યારે $\le$, $+$ અને $\times$
$\Th{PA^{2\dagger}}$ના અનુગામી-માત્ર પાયામાં વ્યાખ્યાયિત કરી શકાય છે.
\end{prop}
\sourcecorrection{OLSOL-006}{સ્થિર અંગ્રેજી મૂળમાં આ ક્રિયાઓ
``સિદ્ધાંતમાં વ્યાખ્યેય'' છે એમ કહેવાયું છે. અર્થ એવો છે કે
અનુગામી અને શૂન્ય ધરાવતા પાયામાં દ્વિતીય-ક્રમ સૂત્રો વડે પ્રમાણભૂત
ક્રમ, સરવાળો અને ગુણાકારના સંબંધો નિર્ધારિત કરી શકાય છે. પ્રથમ બે
સ્વયંસિદ્ધિઓ મનસ્વી રીતે પહેલેથી અર્થઘટિત $+$ કે $\times$ને
નિયંત્રિત કરતી નથી.}

\begin{proof}
$\le$ વ્યાખ્યાયિત છે એમ બતાવવા એવું સૂત્ર~$!A_\le(x, y)$ શોધવું
પડે કે $\Sat{N}{!A_\le(\num n, \num m)}$ ત્યારે અને માત્ર ત્યારે જ
જ્યારે $n \le m$. આ સૂત્ર વિચારો:
\[
!B(x, Y) \ident Y(x) \land \lforall[y][(Y(y) \lif Y(y'))]
\]
સ્પષ્ટ છે કે $!B(\num n, Y)$ કોઈ ગણ $Y \subseteq \Nat$થી ત્યારે
અને માત્ર ત્યારે જ સંતોષાય જ્યારે $\Setabs{m}{n \le m} \subseteq Y$.
તેથી $!A_\le(x, y) \ident
\lforall[Y][(!B(x, Y) \lif Y(y))]$ લઈ શકાય.

સરવાળો વ્યાખ્યાયિત છે એ જોવા નોંધો કે $k+l=m$ ત્યારે અને માત્ર ત્યારે
જ જ્યારે એવું વિધેય~$u$ હોય કે $u(0)=k$, બધા~$n$ માટે
$u(n')=u(n)'$ અને $m = u(l)$. આ સમકક્ષતા વાપરીને
$\Th{PA^{2\dagger}}$માં સરવાળો નીચેના સૂત્ર વડે વ્યાખ્યાયિત કરી
શકાય:
\[
!A_+(x,y,z) \ident \lexists[u][(u(\Obj 0)=x \land \lforall[w][u(w')=u
 (w)'] \land u(y)=z)]
\]
હવે $\Sat{N}{!A_+(\num k, \num l, \num m)}$ ત્યારે અને માત્ર
ત્યારે જ જ્યારે $k+l=m$ એ સ્પષ્ટ હોવું જોઈએ.
\sourcecorrection{OLSOL-007}{સ્થિર અંગ્રેજી મૂળ સરવાળાના સૂત્રમાં
$\lforall[w]$ લખે છે, પરંતુ તેની અંદર પુનરાવર્તન માટે $u(x')=u(x)'$
લખે છે. તેથી $w$ વપરાતો નથી અને જરૂરી સર્વવાચી પુનરાવર્તન મળતું
નથી. અહીં $u(w')=u(w)'$ કર્યું છે.}
\end{proof}

\begin{prob}
\olref[sol][met][spa]{prop:sol-pa-definable}ની સાબિતી પૂરી કરો.
\end{prob}

\end{document}
